\scanpage{42}
dans $\mathfrak g$ sous $G$) ; \deleted{\unsure{de même|donc}, les droites} ``plus généralement, si $g$ est un élément \deleted{régulier principal} de $G$, des \deleted{\unsure{toutes}}\inserted{ses} droites propres\inserted{\uline{régulières}} sous $g$ de $\mathfrak g$ de val. propre une racine prim. $h$.\textsuperscript{ème} existent \deleted{ssi} $g$ est régulier principal, et les droites en question sont conjuguées. On \deleted{\unsure{les}} appelle de telles droites de $\mathfrak g$ cycliques, les pts correspondants ($\ne 0$) de $\mathfrak g$ \uline{cycliques}. \deleted{\unsure{F}} On en construit, avec un couple de Killing donné, en posant $\sum_{1}^{\unsure{r|n}}\lambda_i e_{\alpha_i}+\mu e_{-\psi}$, les $\lambda_i$ et $\mu$ tous $\ne 0$, ($\psi$ la plus haute racine). \unsure{$\Vert$}
% margin note written sideways
\margin{\deleted{\unsure{Correspondant à} tout \unsure{Tore|ss-tore} de $\unsure{G|C}$, \deleted{\unsure{qui est}}\inserted{il} $\exists$\unsure{ra} droite \deleted{\unsure{cyclique}}\inserted{une} stable par $T$ \unsure{d'où un}}}

Le centralisateur d'un élément cyclique\inserted{$x$} de $\mathfrak g$ est un tore maximal $T$, \deleted{le centralisateur est un groupe $\deleted{\unsure{T^{\prime}_{0}}}T'$ entre $T$ et $N(T)$ (ou plutôt il s'agit de la droite $D$ engendrée par $x$) ;} le normalisateur de $D$ dans $G$ est un groupe $T^{\unsure{\sim|\prime}}$, $T\subset T^{\sim}\deleted{\unsure{G'}}\subset N(T)$, $T^{\sim}$ contient l'élément de Coxeter, donc \uline{n'est pas réduit à} $T$. Il faut \deleted{\unsure{donc}} déterminer les ss-groupes correspondants $\deleted{\unsure{W'}} W^{\sim}$ de $W$ ($\deleted{\unsure{W'}}W^{\sim}\subset C^{*}$ est un groupe cyclique\inserted{d'ordre $h'$ multiple de $h$,}\deleted{\unsure{on peut}} \inserted{\unsure{donc en|donc égal à} $h$, grâce à la propriété maximale de l'élément de Coxeter, de \unsure{$T$}.} \deleted{\unsure{on peut voir que}} il est engendré par l'élément de Coxeter. On voit que si $D$ est une droite critique, alors $\unsure{\{}D/\mu_{h'}\simeq E^{1}$ canoniquement [\unsure{se réduire|se rédui|se vérifie}\illegible{1} au cas quasi-déployé, et dans ce cas on a un élément privilégié cyclique, savoir $\sum_i e_{\alpha_i}+e_{-\psi}$. Donc on peut regarder les él\textsuperscript{ts} cycliques qui donnent $1$ dans $\deleted{\unsure{K\simeq|E}}\mathrm{EC}/G\simeq E^{1}$, on peut les appeler les éléments \uline{cycliques principaux} (dépendent à priori du choix d'un signe, \deleted{si $x$ est principal, alors $\unsure{\zeta}x$ est \unsure{antiprincipal}, où $\unsure{\zeta}$ est une racine primitive $2h$.\textsuperscript{ème} de $1$}). Les él\textsuperscript{mts} forment un torseur sous $G$, de stabilisateur un tore maximal $T'$. \deleted{Dans le groupe déployé, \unsure{cet}}\inserted{Pour \deleted{\unsure{que}} un} élément de la forme $\sum e_{\alpha_i}-e_{-\psi}$, on trouve
\margin{\uline{NB} \deleted{\unsure{Le}} le centralisateur d'un sous-tore contenant un élément régulier est un tore maximal}
\margin{\textcircled{?}}
\margin{Mais peut-on choisir un gén. can. de $e_{-\psi}$ ? ici le gpe. des autom \deleted{\unsure{de ce fibré}} d'un gpe épinglé (respectant l'épinglage --) induit-il l'identité sur $P_{-\psi}$ ? \deleted{\illegible{1}} ou mieux sur $P_{\pm\psi}$ ? Sinon, peut-il \ldots}
\margin{préciser comment ?}

e) Considérons les couples $(T,\unsure{g|\gamma})$ d'un tore maximal de $\unsure{G,|G/}$ et d'un élément $g\in N(T)$ \deleted{dont} l'image dans $W$ est un élément de Coxeter. Deux tels couples sont conjugués [\deleted{\unsure{donc}} si on a $(T,g')$ un autre tel couple, $\exists\,n\in N(T)$ tel que $ngn^{-1}=g'$]. Le sous-groupe de $G$ qui stabilise un couple $(T,g)$ est \deleted{discret} un sous-groupe \uline{discret} de $N(T)$, contenant $g$, c'est $N^{\unsure{\sigma|\gamma}}$, extension de $W^{\sigma}$ (car $\mathrm H^1(\mathbb Z/h\mathbb Z,T)=0$) par $T^{\sigma}\deleted{\unsure{est}}\simeq P/(1-\sigma)P\,(1)$ \unsure{--- twist Tate}.
% archive sheet number: 41

